% Part: lambda-calculus
% Chapter: lambda-definability
% Section: lambda-definable-recursive

\documentclass[../../../include/open-logic-section]{subfiles}

\begin{document}

\olfileid{lam}{dfl}{ldr}
\olsection{\usetoken{S}{lambda definable} फलने पुनरावर्ती असतात}

सर्व आंशिक पुनरावर्ती फलने !!{lambda definable}
आहेत एवढेच नव्हे; याचे व्यस्त विधानही खरे आहे.
म्हणजे सर्व !!{lambda definable} फलने आंशिक
पुनरावर्ती असतात.

\begin{thm}
  \ollabel{thm:lambda-computable} आंशिक फलन $f$
  हे !!{lambda definable} असेल, तर ते आंशिक
  पुनरावर्ती असते.
\end{thm}

\begin{proof}
  पुराव्याची फक्त रूपरेषा देऊ. आधी $\lambd$-पदांचे
  अंकगणितीकरण करू: अनुक्रमांच्या नेहमीच्या अभाज्य
  संख्यांच्या घातांवरील संकेतनाचा वापर करून
  $\lambd$-पदांना पद्धतशीरपणे गोडेल संख्या देऊ.
  मग आंशिक पुनरावर्ती फलन $\fn{normalize}(t)$
  परिभाषित करू. ते लॅम्डा पदाची गोडेल संख्या~$t$
  युक्तिवाद म्हणून घेते आणि पदाला सामान्य रूप
  असेल, तर त्या रूपाची गोडेल संख्या परत करते;
  अन्यथा ते अपरिभाषित राहते. पुढे
  $\fn{toChurch}$ आणि $\fn{fromChurch}$ ही दोन
  आंशिक पुनरावर्ती फलने परिभाषित करू. ती नैसर्गिक
  संख्या आणि त्या संख्यांच्या अनुरूप चर्च अंकांच्या
  गोडेल संख्या यांमध्ये पुढे व मागे रूपांतर करतात.

  या पुनरावर्ती फलनांचा वापर करून फलन~$f$
  आंशिक पुनरावर्ती म्हणून परिभाषित करता येते.
  $f$ ला !!{lambda define}s असे $\lambd$-पद~$F$
  आहे. $f(n_1, \dots, n_k)$ संगणित करण्यासाठी
  आधी $\fn{toChurch}(n_i)$ वापरून अनुरूप चर्च
  अंकांच्या गोडेल संख्या मिळवा. त्या $\Gn{F}$
  ला जोडून $F \num{n_1}\dots\num{n_k}$ या पदाची
  गोडेल संख्या मिळवा. आता या गोडेल संख्येवर
  $\fn{normalize}$ वापरा. $f(n_1, \dots, n_k)$
  परिभाषित असेल, तर
  $F \num{n_1}\dots\num{n_k}$ ला सामान्य रूप
  असते (ते चर्च अंकच असले पाहिजे); अन्यथा सामान्य
  रूप नसते (आणि म्हणून
  \[\fn{normalize}(\Gn{F\num{n_1}\dots\num{n_k}})\]
  हे अपरिभाषित असते). शेवटी सामान्यीकृत पदाच्या
  गोडेल संख्येवर $\fn{fromChurch}$ वापरा.
\end{proof}

\end{document}
